\begin{abstracteng}
Climate change has significantly impacted the agricultural sector, particularly in Aceh Besar Regency as one of the major food-producing regions in Aceh Province. Fluctuations in rice production require climate data management that is accurate, consistent, and automated. This study aims to design and implement an automated climate data preprocessing system using BMKG and NASA POWER datasets from 2005\--2025, consisting of 7,565 BMKG records and 7,578 NASA POWER records. The preprocessing stage combined imputation and data smoothing methods tailored to the characteristics of each climate parameter, including Cubic Spline for temperature and humidity, Two-Stage Binary Conditional for rainfall data in the BMKG dataset, and Adaptive Exponential Smoothing to reduce random fluctuations in NASA POWER data. The preprocessing results demonstrated that the applied methods were able to preserve data patterns and trends consistently. In the BMKG dataset, temperature and humidity exhibited excellent smoothing quality with trend preservation above 99\%, while rainfall achieved 90.12\% trend preservation. In the NASA POWER dataset, temperature showed good smoothing quality with 83.5\% trend preservation, whereas humidity maintained trend preservation above 82\% despite exhibiting lower smoothing quality. The results indicate that the developed system is capable of automating climate data preprocessing and producing datasets that are suitable for forecasting applications to support the development of a digital planting calendar.

\bigskip
\noindent
\textit{\textbf{Keywords:}} automated data processing, \textit{smoothing}, GCV, trend preservation, BMKG, NASA POWER
\end{abstracteng}